\section{Exact certification and completion}\label{sec:finite-certification}
\subsection{From theorems to a verifier}
A certificate is a finite witness to one of the preceding necessary
conditions. It is not an unexplained exclusion flag. The basic records are
an irreducible factor and failed divisibility, a prohibited based subring,
a fully checked family recognition, or a finite-field factorization. A
signed record identifies its exhaustive search or shared-cover method,
contains every table remaining after its verified necessary prefilters,
and gives a proof for every remaining character component. An induction record contains all roots of
the codegree polynomial in the dimension field and all admissible low
restriction choices, with a violated Gram or moment condition for each.

For character computations, take
$H=1+b_1+kb_2+k^2b_3$. A separating choice exists for $k\in\{1,\ldots,13\}$:
for each of the six pairs of distinct characters, equality on $H$ excludes
at most two values of $k$. The squarefree characteristic polynomial of
$N_H$ then identifies the rational algebra with
$\Q[t]/(\operatorname{charpoly}(N_H))$. The checker expresses every $N_i$
as a polynomial in $N_H$ and verifies these identities exactly. Factoring
the defining polynomial enumerates all character-field components. The
strictly positive matrix $N_H$ identifies the Frobenius--Perron component
by its largest real eigenvalue, isolated over $\Q$.

There are two different finite completeness checks. The signed search
must cover every bounded coefficient choice modulo its stated gauge; the
low-induction search must cover every vector in
\eqref{eq:column-search} and every allowable column combination. Both are
regenerated during verification. In particular a supplied list of solutions
is never presumed exhaustive. Polynomial roots are checked by multiplying
the linear factors back to the original polynomial. Signs of algebraic
numbers are decided by rational isolating intervals and refinement, not by
floating-point tolerances. Details of the arithmetic implementation and
its elimination branches are recorded in the verification guide.

\begin{proposition}[Soundness of the certificate scheme]\label{prop:certificate-sound}
If the verifier accepts a nonexistence certificate for a fusion ring $R$,
then $R$ has no complex fusion categorification, assuming the established
inputs in Proposition~\ref{prop:standard} and the indicator and central
induction theorems cited in the proofs.
\end{proposition}
\begin{proof}
The direct certificates apply those inputs, the scalar arithmetic lemmas,
Theorem~\ref{thm:nonselfdual-quadratic}, Proposition~\ref{prop:larson-automatic},
or the uniform applications in
Section~\ref{sec:applications}. A weakly integral or two-dimension record
first reduces general existence to pseudo-unitarity, with Galois conjugation
when required. A pseudo-unitary record
applies Proposition~\ref{prop:low-certificate} or the root-length
obstruction, or the global codegree moment bound. A spherical record exhausts all character components,
rejecting zero values, unequal values on duals, a failed norm or field
condition, a Galois orbit with a verified pseudo-unitary obstruction, or
Proposition~\ref{prop:spherical-codegree-moment} at a real embedding.

For a general signed record, Proposition~\ref{prop:signed} supplies at
least one table in the exhaustive bounded search. For a shared search,
Proposition~\ref{prop:global-signed-cover} supplies its global orbit;
the checker verifies that cover before accepting any record depending on
it and tests every compatible transform not already excluded by a common
arithmetic obstruction. A rejected Casimir
factor cannot be its global dimension by Lemma~\ref{lem:norm}; the
prime determinant shortcut is a consequence of the same argument.
The same established $d$-number and cyclotomicity conditions also apply
to the signed quotient, by the joint necessary conditions proved above.
An unsigned gauge yields a spherical categorification, excluded by its
separate spherical record. Otherwise every possible dimension-character
component is rejected by a norm or field certificate or by
Corollary~\ref{cor:sign-certificate}; a whole signed table may instead
be rejected by Corollary~\ref{cor:signed-parity}. Parity is imposed
on the same indicator assignments used in the moment identities, and its
coefficient-domain restriction preserves every categorical possibility.
Thus the hypothetical category has
no remaining branch. These implications have no cyclic dependence on
the classification conclusion or on the companion paper.
\end{proof}

\subsection{The exhaustive run}
\begin{proof}[Proof of Theorem~\ref{thm:finite}]
The complete pinned census \cite{Census} gives
$\CensusSize$ pairwise nonisomorphic based rings of multiplicity at most
$\MultiplicityBound$. The checker reconstructs every full multiplication
tensor and checks the fusion axioms. For each of the fifteen positive entries, the
verifier checks a based-ring isomorphism to the full tensor of a cited
model in Table~\ref{tab:models}. Those established realizations prove
existence.

For every other entry the accepted certificate has one of the types in
Table~\ref{tab:proof-counts}. The counts sum to $\ExcludedCount$.
The $\SignedRecords$ signed records collectively enumerate
$\SignedTables$ normalized associative candidates after their stated
necessary prefilters, with every zero-denominator branch retained.
Candidates omitted before enumeration are eliminated by the stated
necessary filters; every enumerated candidate has a checked rejection. In particular,
neither a ring identifier nor a preassigned classification status is
accepted as a mathematical witness. Proposition~\ref{prop:certificate-sound}
therefore gives general nonexistence for all remaining rings. Unitary
existence for the positive rings and general nonexistence for all others
identify the five existence classifications.
\end{proof}

% Generated from accepted certificates; disjoint first successful tests.
\begin{table}[ht]
\centering
\begin{tabular}{p{.66\linewidth}rr}
\toprule
First successful top-level test, in dispatch order & First use $m$ & Rings\\
\midrule
Formal-codegree $d$-number failure & $1$ & $872\,625$\\
Rank-two moment obstruction & $2$ & $1\,273$\\
Non-self-dual quadratic moments & $2$ & $127$\\
Diagonal involution moments & $3$ & $126$\\
$C_3$ near-group root length & $9$ & $37$\\
Automatic near-group classification & $18$ & $3$\\
Cyclotomicity obstruction & $9$ & $84$\\
Fusion-subring dimension divisibility & $2$ & $28$\\
Residual-involution propagation & $5$ & $1$\\
Weak integrality and induction & $2$ & $1$\\
Two-dimension Galois reduction and induction & $3$ & $8\,126$\\
Signed sphericalization & $2$ & $7\,084$\\
\midrule
Total excluded & & $\ExcludedCount$\\
\bottomrule
\end{tabular}
\caption{Archived disjoint certificate allocation, displayed in the order
of the revised dispatcher. A model is recognized before these tests. First
use is the least exact multiplicity with a certificate of this type;
nested tests in signed records are not counted again. The sequence puts
high-yield direct obstructions before the expensive signed searches.}
\label{tab:proof-counts}
\end{table}
\noindent The finer ledger groups individual checks by their underlying
mathematical criterion. Its latest first use is at multiplicity 18.
In this fixed certificate allocation, all multiplicities from 19 to
$\MultiplicityBound$ require no additional criterion family. This describes
the completed run, not a proof that the collection of criteria is minimal.
For example, signed cyclotomicity is first decisive at multiplicity $95$,
but it applies the same cyclotomicity theorem used earlier for unsigned
characters. Signed-table parity is also used as an enumeration prefilter;
its removed candidate counts are recorded separately from the ring totals.


The numerical tables in this paper are generated from accepted
certificates and the pinned census. The complete per-ring ledger,
per-multiplicity counts, witness files, and criterion-use statistics are
supplied electronically. No manually calculated census case is needed
for the proof. The categorical arguments precede their finite application;
the code only recognizes their hypotheses and checks their exact algebraic
consequences.

\subsection{Compact certificates and reproducibility}
The supplement shares a complete bounded signed census across the
unsigned rings. Proposition~\ref{prop:global-signed-cover} proves that
enumeration modulo all basis-sign changes and nonunit basis permutations
covers every possible signed quotient. Universally invalid signed orbits
are rejected once by formal-codegree $d$-numbers or cyclotomicity.
For each unsigned ring, the verifier expands every remaining orbit into
all compatible transforms and imposes the signed coefficient bounds,
parities, adjoint congruences, and determinant-prime conditions. Every
remaining table receives its own character or induction certificate.
No choice of a convenient orientation substitutes for this full expansion.

The global cover and all its common arithmetic rejections are recomputed
from source during verification; a stored list is not presumed exhaustive.
The smaller direct backends use complete per-ring enumeration, optionally
restricting coefficient domains by adjoint parity first. Every record
specifies its method. In all methods the checker reconstructs the candidate
set, compares it with the complete supplied list, and verifies each longer
certificate. A missing, extra, or altered candidate causes rejection.
The reported candidate count sums the explicitly visited tables after
each method's stated prefilters; it does not count every raw signed
coefficient choice, or count universally excluded orbits repeatedly.
The verified common cover contains $\GlobalSignedOrbits$ signed orbits,
of which $\GlobalSignedSurvivors$ survive the universal arithmetic tests.
The global orbit counts and the method-specific totals are recorded
separately in the verification guide and electronic reports. An independent
arbitrary-integer Python implementation supports comparison with the
direct bounded searches.

Formal-codegree arithmetic is accelerated by exact integer matrix and
polynomial operations. The verifier reconstructs the Casimir polynomial,
checks its irreducible factors and the failed divisibilities, and uses
rational isolating intervals for ordered algebraic comparisons. Floating
point approximations never decide a categorification status. Source
hashes, dependency versions, the pinned census commit, and the commands
for a complete replay are recorded in the verification guide. The
verification of a finite certificate and the census completeness result
documented in \cite{Census} are separate inputs, and both are required.

\begin{conjecture}[Unrestricted classification]\label{conj:classification}
The fifteen rings in Table~\ref{tab:models} are all the rank-four
Grothendieck rings of complex fusion categories, without any multiplicity
restriction.
\end{conjecture}
\begin{conjecture}[Sufficiency of the certificate criteria]\label{conj:criteria}
Every rank-four fusion ring outside Table~\ref{tab:models} has a finite
nonexistence certificate using the necessary criteria and signed reduction
of this paper.
\end{conjecture}
The second conjecture is stronger: it predicts that no additional
categorical obstruction will be needed, with arbitrary exact witnesses
allowed for the stated arithmetic tests. It concerns the mathematical
criteria, not fixed prime-search limits or machine-integer bounds in an
implementation. Its evidence is the complete certificate coverage of
$\ExcludedCount$ exclusions through $\MultiplicityBound$, together with
the unrestricted families handled by the uniform arguments. The two
conjectures assert no theorem above this bound. A uniform obstruction or
bound for the remaining families is still required. The companion paper
develops unrestricted reductions and isolates the unresolved cubic and
quartic pivotal cases.

The supplementary command \texttt{python verify.py} performs the complete
bounded verification. Its default signed-table backend is the supplied C++17 enumerator;
an arbitrary-integer Python backend is provided for independent comparison. The mathematical proof does
not require numerical eigenvalues, pentagon solving or modular-data
classification. The computation certifies the finite algebraic
applications; it is not a proof-assistant formalization of the categorical
arguments.
